% ===========================================================================
%  evnx 0.5.0 — Use Cases
%  Build:  pdflatex -interaction=nonstopmode main.tex   (twice for TOC)
% ===========================================================================
\documentclass[aspectratio=169,10pt]{beamer}

\newcommand{\capturedir}{../captures}
% ---------------------------------------------------------------------------
%  evnx presentation preamble — shared by all three decks
%
%  Engine: pdfLaTeX (also builds under XeLaTeX / LuaLaTeX).
%  Packages: beamer, listings, xcolor, amssymb, newunicodechar, booktabs,
%            fancyvrb, hyperref — all in texlive-latex-recommended + beamer.
%
%  \capturedir must be defined by the including document BEFORE \input of
%  this file, e.g.  \newcommand{\capturedir}{../captures}
% ---------------------------------------------------------------------------

\usetheme{Madrid}
\usecolortheme{seahorse}
\usefonttheme{professionalfonts}
\setbeamertemplate{navigation symbols}{}
\setbeamertemplate{caption}[numbered]
\setbeamerfont{frametitle}{size=\large}
\setbeamerfont{title}{size=\LARGE,series=\bfseries}

\usepackage[T1]{fontenc}
\usepackage[utf8]{inputenc}
\usepackage{lmodern}
\usepackage{amssymb}
\usepackage{booktabs}
\usepackage{xcolor}
\usepackage{listings}
\usepackage{fancyvrb}
\usepackage{newunicodechar}
\usepackage{tikz}
\usetikzlibrary{arrows.meta,positioning,shapes.geometric,fit,backgrounds}

% ---- palette ---------------------------------------------------------------
\definecolor{evnxInk}{HTML}{1F2933}
\definecolor{evnxAccent}{HTML}{2F6F4E}
\definecolor{evnxWarn}{HTML}{B7791F}
\definecolor{evnxErr}{HTML}{B03A2E}
\definecolor{evnxOk}{HTML}{2F6F4E}
\definecolor{evnxMuted}{HTML}{6B7280}
\definecolor{evnxTermBg}{HTML}{F5F5F2}
\definecolor{evnxRule}{HTML}{D1D5DB}

\setbeamercolor{structure}{fg=evnxAccent}
\setbeamercolor{alerted text}{fg=evnxErr}

% ---- glyph mapping for body text ------------------------------------------
% evnx emits exactly these 11 non-ASCII characters (enumerated from the
% captured logs, not guessed).
\newcommand{\evnxCheck}{\textcolor{evnxOk}{\ensuremath{\checkmark}}}
\newcommand{\evnxCross}{\textcolor{evnxErr}{\ensuremath{\times}}}
\newcommand{\evnxWarnSign}{\textcolor{evnxWarn}{\textbf{!}}}

\newunicodechar{✓}{\evnxCheck}
\newunicodechar{✗}{\evnxCross}
\newunicodechar{⚠}{\evnxWarnSign}
\newunicodechar{→}{\ensuremath{\rightarrow}}
\newunicodechar{↔}{\ensuremath{\leftrightarrow}}
\newunicodechar{─}{\textendash}
\newunicodechar{—}{---}
\newunicodechar{·}{\textperiodcentered}
\newunicodechar{…}{\ldots}
\newunicodechar{•}{\textbullet}
\newunicodechar{️}{}                % U+FE0F variation selector — invisible

% ---- terminal listing style ------------------------------------------------
% `literate` is required: newunicodechar does not apply inside listings.
\lstdefinestyle{evnxterm}{
  basicstyle=\ttfamily\scriptsize,
  backgroundcolor=\color{evnxTermBg},
  frame=single,
  rulecolor=\color{evnxRule},
  framesep=4pt,
  breaklines=true,
  breakatwhitespace=false,
  columns=fullflexible,
  keepspaces=true,
  showstringspaces=false,
  upquote=true,
  extendedchars=true,
  literate=
    {✓}{{\textcolor{evnxOk}{$\checkmark$}}}1
    {✗}{{\textcolor{evnxErr}{$\times$}}}1
    {⚠}{{\textcolor{evnxWarn}{\textbf{!}}}}1
    {→}{{$\rightarrow$}}1
    {↔}{{$\leftrightarrow$}}1
    {─}{{-}}1
    {—}{{---}}1
    {·}{{$\cdot$}}1
    {…}{{\ldots}}1
    {•}{{\textbullet}}1
    {️}{{}}0
    {é}{{\'e}}1
}

\lstdefinestyle{evnxcode}{
  style=evnxterm,
  basicstyle=\ttfamily\scriptsize,
  backgroundcolor=\color{white},
}

\lstset{style=evnxterm}

% ---- helper macros ---------------------------------------------------------
% \capture{file}        — include a real captured run verbatim
% \capturepart{f}{a}{b} — include lines a..b of a captured run
\newcommand{\capture}[1]{\lstinputlisting[style=evnxterm]{\capturedir/#1}}
\newcommand{\capturepart}[3]{%
  \lstinputlisting[style=evnxterm,firstline=#2,lastline=#3]{\capturedir/#1}}

% a labelled "real output" box
\newcommand{\realrun}[1]{%
  \begin{center}\scriptsize\textcolor{evnxMuted}{real run \textendash\ #1}\end{center}}

\newcommand{\cmd}[1]{\texttt{\textbf{#1}}}
\newcommand{\flag}[1]{\texttt{#1}}
\newcommand{\file}[1]{\texttt{#1}}

% exit-code chips
\newcommand{\exitzero}{\textcolor{evnxOk}{\texttt{0}}}
\newcommand{\exitone}{\textcolor{evnxWarn}{\texttt{1}}}
\newcommand{\exittwo}{\textcolor{evnxErr}{\texttt{2}}}

% a standard "verified against" footline for factual slides
\newcommand{\verifiedon}{\tiny\textcolor{evnxMuted}{verified against evnx 0.5.0 @ 9cfe75b}}

\usepackage[colorlinks=true,linkcolor=evnxAccent,urlcolor=evnxAccent]{hyperref}


\title{evnx in practice}
\subtitle{Workflows that combine commands --- and what cloud sync adds}
\date{}

\begin{document}

\begin{frame}[plain]
  \titlepage
  \begin{center}\footnotesize
  \textcolor{evnxMuted}{evnx 0.5.0 @ \texttt{9cfe75b}. Every terminal block is a
  real run.}
  \end{center}
\end{frame}

\begin{frame}{Why this deck exists}
  The command guide shows each command alone. Almost nothing useful is one
  command.

  \vspace{0.8em}
  \begin{block}{Five workflows}
    \begin{enumerate}
      \item \textbf{Day one} — a project with no \file{.env} discipline at all
      \item \textbf{The breaking PR} — catching a leak before review
      \item \textbf{The CI gate} — annotations, SARIF, and what fails open
      \item \textbf{Many environments} — staging vs production drift
      \item \textbf{Cloud sync} — a team, a pipeline, and no plaintext file
    \end{enumerate}
  \end{block}

  \vspace{0.6em}
  {\footnotesize Each ends with the exact command chain, copy-pasteable.}
\end{frame}

\begin{frame}{Contents}\tableofcontents[hideallsubsections]\end{frame}

% =========================================================================
\section{Day one}

\begin{frame}{The situation}
  A repository you just cloned. There is a \file{.env} on someone's laptop, a
  \file{.env.example} that is six months stale, and a README that says
  ``ask Priya for the keys''.

  \vspace{0.8em}
  \begin{block}{What you want}
    A \file{.env.example} that is true, a \file{.gitignore} that is correct, and
    a way to tell whether your local \file{.env} is actually complete.
  \end{block}
\end{frame}

\begin{frame}[fragile]{\cmd{init --from-source}: ask the code, not the README}
  \capturepart{80-uc1-onboard.txt}{1}{20}
\end{frame}

\begin{frame}[fragile]{What it produced}
  \capturepart{80-uc1-onboard.txt}{21}{48}

  {\footnotesize Each variable carries the \texttt{file:line} where the code
  reads it. That is the difference between a template someone wrote once and a
  template derived from the program.}
\end{frame}

\begin{frame}[fragile]{The chain}
\begin{lstlisting}[style=evnxcode,language=bash]
git clone <repo> && cd <repo>

evnx doctor                      # gitignore, permissions, structure
evnx init --from-source          # propose a template from the code
evnx init --from-source --yes    # accept it

evnx sync --direction reverse --force   # fill .env from the template
$EDITOR .env                            # put the real values in
evnx validate                           # did I miss any?
\end{lstlisting}

  \begin{block}{Why \flag{--from-source} and not \flag{--detect}}
    \flag{--detect} reads \file{package.json} and infers a stack.
    \flag{--from-source} reads \texttt{process.env.X} in the actual code. Only
    the second can answer ``is my template complete?''.
  \end{block}
\end{frame}

% =========================================================================
\section{The breaking PR}

\begin{frame}[fragile]{The situation}
  Nobody types \texttt{LEAK\_MY\_KEY=true}. Someone pastes a key from the Stripe
  dashboard to debug locally, and flips \texttt{DEBUG} while they are in there
  --- and quotes it, because every other value on the page is quoted.

  \vspace{0.8em}
\begin{semiverbatim}\small
+ STRIPE_SECRET_KEY=sk_live_51QwErTy...
+ DEBUG="False"
\end{semiverbatim}

  \vspace{0.6em}
  {\footnotesize Two lines. Three problems: a live key in a tracked-adjacent
  file, a duplicate assignment that silently wins, and a boolean that is
  \texttt{True} in Python.}
\end{frame}

\begin{frame}[fragile]{Before: green}
  \capture{81-uc2-before.txt}
\end{frame}

\begin{frame}[fragile]{\cmd{scan} catches the leak}
  \capture{82-uc2-scan.txt}

  {\footnotesize Note it distinguishes \texttt{(test)} from \texttt{(LIVE)} and
  names where to revoke. Values are masked.}
\end{frame}

\begin{frame}[fragile]{\cmd{validate} catches the other two}
  \capturepart{83-uc2-validate.txt}{1}{18}
\end{frame}

\begin{frame}[fragile]{The boolean trap, in one line}
\begin{semiverbatim}\small
>>> import os
>>> os.getenv("DEBUG")
'False'
>>> bool(os.getenv("DEBUG"))
True          \# <- debug mode is ON in production
\end{semiverbatim}

  \vspace{0.6em}
  \texttt{bool} of a non-empty string is \texttt{True}. Always. \texttt{'False'}
  is a non-empty string.

  \vspace{0.8em}
  \begin{block}{And the duplicate}
    \texttt{STRIPE\_SECRET\_KEY is assigned 2 times (lines 3, 5) — only line 5
    takes effect.} In a 60-line file that is invisible, and the value someone
    just carefully edited is the one that loses.
  \end{block}
\end{frame}

\begin{frame}[fragile]{The chain: a pre-commit hook}
\begin{lstlisting}[style=evnxcode,language=bash]
cat > .git/hooks/pre-commit <<'EOF'
#!/bin/sh
evnx scan --severity high || {
  echo "evnx: secrets detected - commit aborted"; exit 1; }
EOF
chmod +x .git/hooks/pre-commit
\end{lstlisting}

  Or through \texttt{pre-commit}:

\begin{lstlisting}[style=evnxcode,language=yaml]
repos:
  - repo: https://github.com/urwithajit9/evnx
    rev: v0.5.0
    hooks:
      - id: evnx-scan
      - id: evnx-validate
\end{lstlisting}
\end{frame}

% =========================================================================
\section{The CI gate}

\begin{frame}[fragile]{A workflow that blocks the merge}
\begin{lstlisting}[style=evnxcode,language=yaml]
name: env-guard
on: pull_request
jobs:
  guard:
    runs-on: ubuntu-latest
    permissions:
      contents: read
      security-events: write        # required for the SARIF upload
    steps:
      - uses: actions/checkout@v4
      - run: curl -fsSL https://dotenv.space/install.sh | sh

      - run: evnx scan --severity high --format github    # blocks the PR

      - if: always()
        run: evnx scan --format sarif --exit-zero > evnx.sarif
      - if: always()
        uses: github/codeql-action/upload-sarif@v3
        with: { sarif_file: evnx.sarif }

      - run: evnx sync --check
\end{lstlisting}
\end{frame}

\begin{frame}[fragile]{The same steps, run for real}
  \capture{85-uc3-ci.txt}
\end{frame}

\begin{frame}{Three things that cost real debugging time}
  \begin{enumerate}
    \item \textbf{\flag{--exit-zero} on the SARIF step is not optional.}
          Without it the step fails, a failed step means
          \texttt{upload-sarif} never runs, and your Security tab stays empty
          while you believe it is working.
    \item \textbf{Two scan runs, deliberately.} The first blocks the pull
          request with inline annotations. The second must not fail.
    \item \textbf{\flag{--severity high} is the gate.} It filters what is
          reported, counted \emph{and} exited on. The default (\texttt{low})
          blocks on every high-entropy string — including public keys that are
          meant to be public.
  \end{enumerate}

  \vspace{0.5em}
  \begin{alertblock}{And one that does not exist}
    There is no \flag{--exit-code} and no \flag{--path}. Paths are positional;
    the exit code is the default behaviour and \flag{--exit-zero} is the
    opt-\emph{out}.
  \end{alertblock}
\end{frame}

% =========================================================================
\section{Many environments}

\begin{frame}[fragile]{Drift between staging and production}
  \capture{86-uc4-envdiff.txt}

  {\footnotesize \flag{--against} compares two \emph{real} environments and
  exits \exitzero{} when they match — unlike the default comparison against a
  template, where every filled-in value differs by construction.}
\end{frame}

\begin{frame}[fragile]{Context-aware checks}
  \capturepart{87-uc4-prodcheck.txt}{1}{12}

  {\footnotesize This fires only because a \file{docker-compose.yml} is present.
  \texttt{localhost} is correct on a laptop and wrong inside a container, so the
  same value is a finding in one project and not in another.}
\end{frame}

\begin{frame}[fragile]{The chain}
\begin{lstlisting}[style=evnxcode,language=bash]
# is production missing anything staging has?
evnx diff --env-name production --against staging \
     --ignore-keys APP_ENV,FEATURE_NEW_CHECKOUT

# is each environment internally valid?
for e in development staging production; do
  evnx validate --env-name "$e" --format json 2>/dev/null \
    | jq -r '.issues[].message'
done

# render the real config for one of them
evnx template --input app.yaml.tpl --output app.yaml \
     --env-name production --strict
\end{lstlisting}
\end{frame}

% =========================================================================
\section{Cloud sync}

\begin{frame}{What the cloud adds --- and what it does not}
  \begin{columns}[T]
    \begin{column}{0.48\textwidth}
      \textbf{It adds}
      \begin{itemize}\footnotesize
        \item One command to get the team's \file{.env}
        \item Versioned history you can roll back
        \item Sharing that can be revoked
        \item Secrets into CI with no file on disk
      \end{itemize}
    \end{column}
    \begin{column}{0.48\textwidth}
      \textbf{It does not add}
      \begin{itemize}\footnotesize
        \item Trust in the server --- it holds ciphertext
        \item Freedom from the master password
        \item Protection from a malicious server \emph{when sharing}
      \end{itemize}
    \end{column}
  \end{columns}

  \vspace{0.8em}
  \begin{block}{The guarantee}
    Encryption happens on your machine. The server stores opaque blobs and
    cannot read your secrets --- not with full database access.
  \end{block}
\end{frame}

\begin{frame}[fragile]{The chain: a team}
\begin{lstlisting}[style=evnxcode,language=bash]
evnx auth register                       # keys generated locally
evnx auth totp enable                    # QR in the terminal
evnx vault create shop/production
evnx cloud link shop/production          # bind this directory

evnx cloud push                          # encrypt .env, upload ciphertext
evnx vault share shop/production --user teammate@example.com

# teammate, on another machine
evnx auth login && evnx cloud pull       # .env.production appears
\end{lstlisting}

  \begin{block}{Naming}
    A vault is \texttt{name} or \texttt{name/environment}. Pushing
    \texttt{shop/production} picks up \file{.env.production}, and cloud commands
    always print which file they chose.
  \end{block}
\end{frame}

\begin{frame}[fragile]{The chain: a pipeline, with no plaintext file}
\begin{lstlisting}[style=evnxcode,language=bash]
# a read-only token scoped to one vault
evnx auth token create --scope read_only --vault shop/production

# in CI - the secrets never touch the filesystem
echo "$EVNX_MASTER_PASSWORD" | \
  evnx cloud run --vault shop/production --password-stdin -- ./deploy.sh

# least privilege: the build step sees only the public variables
echo "$EVNX_MASTER_PASSWORD" | \
  evnx cloud run --vault shop/production --password-stdin \
       --include 'NEXT_PUBLIC_*' -- npm run build
\end{lstlisting}

  \vspace{0.3em}
  {\footnotesize Values travel in the environment block --- not \texttt{argv} ---
  so nothing appears in \texttt{ps}, in shell history, or on disk. A filter
  matching nothing is an \textbf{error}, not a silent run with zero secrets.}
\end{frame}

\begin{frame}[fragile]{Verify the claim yourself}
\begin{lstlisting}[style=evnxcode,language=bash]
# the secret must NOT appear in the child's command line
evnx cloud run --vault shop/production -- \
     sh -c "tr '\\0' '\\n' < /proc/self/cmdline"

# but it IS in the environment the child was handed
evnx cloud run --vault shop/production -- printenv DATABASE_URL
\end{lstlisting}

  \vspace{0.4em}
  Published recipe:
  \url{https://www.evnx.dev/guides/reference/verifying-the-round-trip}

  \vspace{0.6em}
  \begin{alertblock}{Say the limit out loud}
    \cmd{cloud run} removes the plaintext \emph{file}, not the master password.
    Anyone who tells you their secrets tool eliminated the root secret is
    selling you something.
  \end{alertblock}
\end{frame}

\begin{frame}[fragile]{The transport guard}
  \capture{66-tls-guard.txt}

  {\footnotesize Loopback is exempt so local development works. The error
  explains the \emph{threat}, not just the rule.}
\end{frame}

\begin{frame}{The honest caveat on sharing}
  \begin{alertblock}{The recipient's public keys come from the server}
    A malicious server could substitute its own and read what you share. There
    is no third party to check against, and out-of-band fingerprint verification
    is not built.
  \end{alertblock}

  \vspace{0.6em}
  So \emph{``the server cannot read your secrets''} becomes \emph{``\ldots cannot
  read them \textbf{passively}''} the moment you share.

  \vspace{0.8em}
  \begin{block}{What is solid}
    A solo vault is wrapped under your master key --- symmetric only, so
    post-quantum safe. A shared vault uses hybrid X25519 + ML-KEM-768 with
    \textbf{no} classical-only fallback, and revoking a member re-keys the
    vault so they cannot read versions pushed after removal.
  \end{block}
\end{frame}

% =========================================================================
\section*{Closing}
\begin{frame}{Pick the workflow that fits}
  \begin{tabular}{@{}p{0.34\textwidth}p{0.60\textwidth}@{}}
    \toprule
    \textbf{You are} & \textbf{Start with} \\
    \midrule
    Joining a project      & \texttt{evnx doctor \&\& evnx init --from-source} \\
    Reviewing a PR         & \texttt{evnx scan --severity high --format github} \\
    Wiring CI              & \texttt{scan --format sarif} + \texttt{sync --check} \\
    Managing environments  & \texttt{evnx diff --env-name X --against Y} \\
    Sharing with a team    & \texttt{evnx vault share} + \texttt{cloud push/pull} \\
    Deploying              & \texttt{evnx cloud run -- ./deploy.sh} \\
    \bottomrule
  \end{tabular}

  \vspace{1em}
  \begin{center}
  \url{https://www.evnx.dev} \quad\textperiodcentered\quad
  \texttt{curl -fsSL https://dotenv.space/install.sh | sh}
  \end{center}
\end{frame}

\end{document}
